\begin{center}
{\Large\bfseries 摘要}
\end{center}

随着宏基因组学技术的发展，研究人员已经能从未经培养的人类相关样本中准确估计微生物群落。这为人类疾病预测和生物标志物发现提供了丰富的微生物特征。本文基于结直肠癌（CRC）、炎症性肠病（IBD）与肥胖症（Obesity）三个独立宏基因组，共484个样本，1331个物种级特征，研究疾病在菌群丰度特征下的可预测性、每类疾病的特异标志菌群特征，以及预测模型的跨疾病泛化能力。题目数据呈高维稀疏、样本量远小于特征数且相对丰度存在定和约束，直接建模易致过拟合与伪相关。

针对以上挑战，我们先剔除了零值占比大于95\%的特征，将指标从1331维降低到264维，接着通过中心化对数比变换（CLR）破除定和约束，消除丰度数据的伪相关性，然后逐步对三个问题建模。
问题一以L2正则 Logistic 回归为主模型，以随机森林为对照，选取AUC为主指标评估疾病预测性能。结直肠癌、炎症性肠病、肥胖症的 AUC 分别达到 0.7907、0.8871 与 0.6496（RF最优达 0.9035）。结果经过交叉验证与LOOCV检验，两者结果自洽，无过拟合。
问题二采用 $L_1$ 正则化结合 Bootstrap 稳定性选择筛选标志物，通过 Fisher 精确检验存在/缺失差异；用Wilcoxon秩和检验丰度差异，并通过BH-FDR 校正控制多重比较的假阳性率。我们与文献报道的疾病相关菌比对后发现结直肠癌和炎症性肠病分别命中 3/4 和 4/4 的已知菌种。该方法在生物学领域有效。
问题三规定留一疾病（leave-one-disease-out, LODO）框架，评估泛化能力。跨疾病通用模型最优 AUC 仅 0.607，低于可用性参考线 0.65；经六轮归因实验排除了批次与算法表征等外在干扰，最终证实泛化性能不足来自疾病信号高度特异以及数据固有的限制。

本文通过多维扰动实验对结论做复核。在生物标志物筛选中，入选特征对选择阈值 $\tau\in[0.4,0.7]$ 高度稳健，CRC 与 IBD 的核心标志物重合度达 100\%，并且在 PLS-DA 的 $\mathrm{VIP}>1.5$ 排序中得到一致佐证。跨疾病预测在尝试树模型、样本合并与密度比加权等方法后，最优 AUC 稳定收敛于 0.6068，验证了数据限制。

综上，本文研究揭示了宏基因组的疾病判别信号主要由物种的“存在/缺失”模式主导，且三病入选的稳定标志物集合间 Jaccard 相似系数为 0，失调菌群呈现高度的疾病特异性。上述发现从机理上决定了跨疾病预测模型的泛化能力受限。本文方法兼顾了生物可解释性与建模自洽性，为临床单病种宏基因组筛查提供了坚实的方法论依据。

\vspace{0.6em}
\noindent{\bfseries 关键词：}宏基因组；疾病预测；生物标志物；CLR 变换；稳定性选择；跨疾病泛化
